\section{Failure Attribution}
\label{sec:failure-attribution}

The taxonomy in Section~\ref{sec:multi-agent-trace} names a failure mode but does not indicate which agent or which step in the trace caused it. Failure attribution is the additional task of assigning a trace failure to the agent and step responsible for it, so that a fix can target that agent or step rather than the run as a whole.

\subsection{Attribution Methods}
\label{sec:attribution-methods}

Zhang et al. formulated this step and introduced Who\&When, a benchmark of 127 annotated multi-agent systems in which each failed trace is labeled with the agent and the step responsible for the failure \cite{zhangWhichAgentCauses2025}. The automated methods that they evaluate on this benchmark are themselves built on the LLM-as-a-judge paradigm of Section~\ref{sec:llm-as-judge}, applied here to read a failure trace and return a verdict on who is responsible instead of grading a single output \cite{zhangWhichAgentCauses2025}. Against this benchmark, the best of these methods reached 53.5\% accuracy at identifying the responsible agent and 14.2\% accuracy at pinpointing the responsible step, and some methods performed below random on the step-level task \cite{zhangWhichAgentCauses2025}. The gap between the two figures indicates that identifying which agent is at fault is considerably easier than identifying at which point in that agent's trajectory the fault occurs. The low step-level accuracy is a further instance of the judge unreliability already established in Section~\ref{sec:llm-as-judge}, here applied to a harder task, reading a failure trace instead of grading a single response. These methods are trained and evaluated against the 127 traces in Who\&When. Labeling a trace at this granularity (agent and step) by hand is costly, which limits how much attribution data of this kind can exist.

AgenTracer targets this labeling cost: it automates the annotation of failed multi-agent trajectories through counterfactual replay and programmed fault injection, rather than relying on traces labeled by hand \cite{zhangAgenTracerWhoInducing2025}. Fault injection here means taking a trajectory that succeeded and deliberately corrupting one of its actions, so that the resulting failure has a known cause by construction \cite{zhangAgenTracerWhoInducing2025}. A tracer trained on the traces that the AgenTracer pipeline produces outperformed large proprietary judges, including Gemini-2.5-Pro and Claude-4-Sonnet, by up to 18.18\% on the Who\&When benchmark \cite{zhangAgenTracerWhoInducing2025}.

Both Who\&When's judges and AgenTracer's automated labels treat attribution as the selection of one responsible agent and one responsible step, a single ground truth to predict. DoVer questioned this framing and approached attribution as intervention-driven debugging instead: it generates a hypothesis about the responsible agent and step, then verifies the hypothesis through targeted interventions on the trace, in the same spirit as AgenTracer's fault injection but applied to check a hypothesis rather than to generate training data, such as editing a message or altering a plan, rather than accepting the hypothesis on its own \cite{maDoVerInterventionDrivenAuto2026}. The authors observed that framing attribution as the search for a single responsible agent and step is often ill-posed, since multiple distinct interventions can independently repair the same failed task \cite{maDoVerInterventionDrivenAuto2026}. In et al. reported the same difficulty from a different angle: existing benchmarks and methods largely assume a single deterministic root cause, but multi-agent system failures often admit multiple plausible attributions because of complex inter-agent dependencies and ambiguous execution trajectories \cite{inRethinkingFailureAttribution2026}. A trace failure, in this view, is not always reducible to one agent and one step; more than one attribution can be defensible for the same run.

\subsection{Cost of Attribution}
\label{sec:cost-of-attribution}

Attribution methods of the kind surveyed above need a complete trace to reason over, and reasoning over the trace once is often not enough: both AgenTracer and DoVer establish the responsible agent and step by replaying the trace under a modification rather than by reading it once. AgenTracer produces its labels through counterfactual replay combined with programmed fault injection \cite{zhangAgenTracerWhoInducing2025}, and DoVer replays each trial with an intervention applied in-situ, running three independent replays per intervention and, across the four settings it evaluates, between 1.4 and 4.0 replay trials per case \cite{maDoVerInterventionDrivenAuto2026}. Producing an attribution this way therefore implies two things. First, it requires re-executing part or all of the trace multiple times. Each replay checkpoints the state it starts from, and the procedure repeats until a hypothesis is confirmed. This adds a compute cost on top of whatever was already spent producing the original trace, and the cost scales with the number of interventions or fault injections a method tries before it converges on an attribution. Second, replaying a trace requires the trace to already exist in full: the run must have finished before any replay, checkpoint, or intervention can be performed on it. Attribution of this kind is therefore an offline procedure, carried out after the run is complete, and its cost in money and time is paid only after a failure has already happened. This implies a distinction between offline evaluation, performed on a completed trace, and online evaluation, performed during execution, developed in Section~\ref{sec:offline-online-evaluation}.

